Major depressive disorder (MDD) remains a leading cause of global disability and is characterized by marked heterogeneity in clinical presentation, neurobiology, and treatment response. Despite the availability of pharmacological and psychotherapeutic interventions, a substantial proportion of patients fail to achieve remission, motivating increasing interest in neuromodulation approaches such as transcranial direct current stimulation (tDCS) \citep{Woodham2021}. While randomized controlled trials and recent large-scale studies support the feasibility and clinical potential of tDCS, including fully remote, home-based delivery, treatment outcomes remain variable at the individual level \citep{woodham2025home,Hampsey2026}.

Home-based delivery has moved from concept to clinical trial: a fully remote phase 2 randomized sham-controlled study demonstrated home-based tDCS in major depressive disorder \citep{woodham2025home}, and qualitative follow-up work found the approach acceptable to participants \citep{Rimmer2024}. These advances make pretreatment stratification a practical clinical problem. A central challenge is pretreatment stratification: identifying which patients are more likely to achieve remission before treatment begins. Electroencephalography (EEG) is a potential modality for this purpose due to its portability, low cost, and high temporal resolution.

The literature on EEG biomarkers in major depressive disorder is encouraging though methodologically uneven. Reviews and biomarker overviews have repeatedly highlighted spectral, asymmetry, vigilance, entropy, and complexity signatures as potentially informative marker families \citep{olbrich2013_eeg,olbrich2015_personalized,leuchter2010_biomarkers}. Related discussions of neuroimaging and EEG biomarkers have emphasized the clinical opportunity, but also the need for more robust treatment-response models and better translational discipline \citep{Fu2013,Fu2019,Widge2019}. More recent meta-analytic work on EEG-based treatment prediction reaches a similar conclusion: useful signal is present, but translation depends on stronger methodology and standardization \citep{simmatis2023_technical,watts2022predicting,cohen2023electroencephalography}. Machine learning studies have extended this literature into antidepressant, placebo, and neuromodulation response prediction, while also exposing recurrent concerns about external generalization, feature stability, leakage-safe selection, and interpretability \citep{AlKaysi2017,jaworska2019_leveraging,oakley2022_eeg}. Related nonlinear and entropy-oriented studies point to the same caution about dimensionality and signal instability \citep{shalbaf2018_nonlinear,ebrahimzadeh2021_predicting,zhdanov_svm_prediction}. Entropy-focused EEG work has reinforced the relevance of complexity measures \citep{Puthankattil2012,Ghiasi2021,Zhao2022}. Leakage-specific and data-centric critiques make the underlying design problem explicit \citep{Shim2021,Shen2025}.

That tension is especially sharp in relatively small neuromodulation studies. Rich models can improve representation learning, but relatively small cohorts also increase the risk of unstable feature selection, inflated performance estimates, and overly confident biomarker claims. Nested evaluation and within-fold processing are therefore central design issues rather than implementation details in this literature \citep{Widge2019,Shim2021}. Clinically deployable EEG biomarkers also need interpretability and low-density feasibility, and recent data-centric frameworks have started to move in that direction \citep{Shen2025}. Emerging response-prediction studies in depression and related neuromodulation settings show substantial promise \citep{funk2008_prefrontal,bailey2023_concurrent,murphy2023_neural}. Related reviews and pilot studies also make the methodological fragility clear \citep{guo2024_restingstate,tsai2023_review,Dagnino2023}. More recent response-oriented studies continue that pattern \citep{sheen2024_response,zeidabadi2025_prediction,reissmann2026_thetacordance}. A recent systematic review of EEG connectivity changes in depression treatment reaches a similar conclusion \citep{ulrich2025_alterations}.

Within the same home-based 4-channel MDD program analysed here, power spectral density (PSD)-based and PLV-based analyses have observed that the baseline EEG carries treatment-response information across multiple signal representations \citep{Moncy2024,Xiao2025PMIP}. PSD-based deep learning approaches identified temporoparietal spectral features, notably TP10 theta-alpha power, as informative for treatment outcome prediction \citep{Moncy2024}, while connectivity-based analyses highlighted frontal-temporal and temporoparietal network patterns \citep{Xiao2025PMIP}. Together, these findings suggest that response-related signal is present and reproducible across feature families, but remains poorly characterized in terms of interpretability, stability, and minimal feature representation.

The present study builds directly on this context and uses the same baseline EEG cohort as our prior PSD- and connectivity-based analyses \citep{Moncy2024,Xiao2025PMIP}, but addresses a distinct question. Rather than optimising predictive performance within a predefined feature family, we test whether a strictly leakage-controlled, foldwise sparse feature-selection pipeline can identify a small, recurrent, and interpretable biomarker set under participant-level validation. We compared two classifiers over the same feature-selection sweep to test whether low-montage resting-state EEG could recover a small recurring set of biomarkers associated with tDCS remission in adults with MDD. The analysis prioritized parsimony, biomarker interpretability, and leakage-resistant participant-level validation. Our working hypothesis was that useful patient-level discrimination could be achieved with a small per-fold feature budget and that the most recurrent features would show stronger effect-direction consistency than hard-set overlap. This approach reframes biomarker discovery from a high-dimensional modelling problem toward a minimal, interpretable signal identification problem.
